\section{Empirical robustness and diagnostics}
\label{sec:empirical-robustness}

This appendix first audits the dyadic estimator and inference, then turns to
representation and finite article sampling, and finally considers the sector
projection and exploratory return-path display used in \Cref{sec:results}.
These diagnostics preserve the main estimand and do not add independent theorem tests.

\subsection{Dyadic specification and inference audit}
\label{sec:dyadic-diagnostics}
\label{sec:resampling-formulas}

For the reported coefficient, let $X$ and $D$ stack $X_{ij}$ and $D_{ij}$ over
the observed dyads $i<j$, and let $F$ collect the intercept and one indicator
per non-omitted firm from the rank-cleaned point design. This nuisance design
is distinct from the secondary-control vector $z_{ij}$ in
equation~\ref{eq:dyadic-fixed-effects}, which is zero in the canonical
specification. Write $M_F=I-F(F^{\top}F)^{-}F^{\top}$, with
$\widetilde X=M_F X$ and $\widetilde D=M_F D$. The
Frisch--Waugh--Lovell identity is
\begin{equation}\label{eq:fwl-identity}
  \widehat\beta_X
  =\frac{\widetilde X^{\top}\widetilde D}{\widetilde X^{\top}\widetilde X}.
\end{equation}
In words, both axes are first stripped of their average firm-endpoint
components. The remaining slope therefore measures whether unusually distant
article distributions for a given pair coincide with unusually weak return
co-movement for that pair.

The nuisance block has \ppvalue{confound-fwl-nuisance-columns} columns and
omits the indicator for \texttt{\ppvalue{confound-fwl-omitted-firm}}.

The conditional, or partial, $R^2$ of the article-embedding regressor is
\begin{equation}\label{eq:partial-r2}
  R^2_{\mathrm{partial}}
  =1-\frac{\mathrm{SSE}_{\mathrm{FE}+X}}{\mathrm{SSE}_{\mathrm{FE}}},
\end{equation}
where $\mathrm{SSE}_{\mathrm{FE}}$ is the residual sum of squares from the
intercept and symmetric firm effects, and $\mathrm{SSE}_{\mathrm{FE}+X}$ is
from that model with $X$ added. It is defined only for the restricted
fixed-effect model and is therefore omitted on pooled rungs.

The six-rung W$_2$ ladder asks whether the association between
article-embedding distance and return distance is confined to raw sample
composition or remains after increasingly stringent adjustment for observable
and stable firm-level heterogeneity.
The sequence is increasingly demanding in the heterogeneity it absorbs,
although the final firm-effect specification changes the adjustment design
rather than simply adding another pooled control.
Specification 1 records the unadjusted W$_2$ association.
Specification 2 adds a same-sector indicator to address broad sector sorting;
because sector remains a coarse description of firm composition,
specification 3 introduces the declared pooled controls for firm-size
differences and crisis-window coverage.
The preceding firm and window controls do not address differential news
coverage, so specification 4 adds declared article-count level and imbalance controls.
Article counts do not capture within-firm content dispersion, so
specification 5 adds the sum and absolute difference of that dispersion.
This decomposition diagnostic separates between-firm transport distance from
heterogeneity in the breadth of the two article clouds.
Because the pooled controls address only their measured dimensions,
specification 6 replaces them with symmetric firm effects that absorb stable
additive endpoint heterogeneity.
Only this sixth specification is primary; the preceding five document the
path from raw association to the preferred design.

\Cref{tab:dyadic-ladder} reports these six prespecified specifications.
Pooled and firm-effect $R^2$ values are not comparable because the latter
absorb endpoint margins; partial $R^2$ therefore compares the focal regressor
only with the restricted firm-effect model.

% AUTO-GENERATED by pipeline render-paper1-tables -- do not edit by hand.
\begin{table}[H]
\centering
\small
\begin{tabularx}{\linewidth}{r>{\raggedright\arraybackslash}Xrrrrrr}
\toprule
Spec. & Adjustment & $\hat\beta_{W_2}$ & $\Delta D$/1 SD & 95\% CI & $p$ & $R^2_{\mathrm{full}}$ & $R^2_{\mathrm{partial}}$ \\
\midrule
\multicolumn{8}{l}{\textit{Panel A: staged pooled adjustments}} \\
1 & W2 distance only & 1.746 & 0.064 & [1.302, 2.271] & 0.001 & 0.258 & \textemdash \\
2 & + same-sector indicator & 1.373 & 0.050 & [0.964, 1.853] & 0.001 & 0.288 & \textemdash \\
3 & pooled controls & 1.461 & 0.054 & [1.074, 1.953] & 0.001 & 0.298 & \textemdash \\
4 & + article-count controls & 1.461 & 0.054 & [1.074, 1.953] & 0.001 & 0.298 & \textemdash \\
5 & + within-firm dispersion ($B_i+B_j$, $|B_i-B_j|$) & 1.782 & 0.065 & [1.351, 2.244] & 0.001 & 0.324 & \textemdash \\
\midrule
\multicolumn{8}{l}{\textit{Panel B: canonical endpoint-adjusted specification}} \\
6 & W2 metric + symmetric firm effects & 1.874 & 0.069 & [1.516, 2.250] & 0.001 & 0.742 & 0.332 \\
\bottomrule
\end{tabularx}
\caption{Dyadic association ladder. Every interval and two-sided sign-tail probability uses the same deterministic multinomial node-count array across rungs. Panel A separates the W2 metric's staged pooled adjustments from canonical specification 6, which absorbs symmetric additive firm effects. Specification 5 conditions on the within-firm dispersion components of article geometry and is a decomposition diagnostic. $\Delta D$/1 SD is the fitted change in return chord distance for a one-standard-deviation change in rooted W2 distance. $R^2_{\mathrm{full}}$ is the full-model coefficient of determination; the pooled and fixed-effect blocks absorb different variation and are compared within block. $R^2_{\mathrm{partial}} = 1 - \mathrm{SSE}_{\mathrm{FE}+W_2}/\mathrm{SSE}_{\mathrm{FE}}$ is the conditional fit against the restricted firm-effect model and is undefined (\textemdash) for the pooled rungs, which have no such restricted model.}
\label{tab:dyadic-ladder}
\end{table}


The focal coefficient remains positive throughout, showing that none of the
declared adjustments reverses the sample association.
This sign stability does not establish causality or rule out pair-specific confounding.

To distinguish substantive specification changes from rank or conditioning
problems, \Cref{tab:dyadic-design-diagnostics} records the observation count,
design rank, scaled conditioning, and residualized W$_2$ variation for every
specification.
The rank-cleaned canonical design omits one firm indicator, and the
Frisch--Waugh--Lovell coefficient agrees with the direct least-squares fit.

% AUTO-GENERATED by pipeline render-paper1-tables -- do not edit by hand.
\begin{table}[H]
\centering
\footnotesize
\begin{tabularx}{\linewidth}{@{}lX@{}}
\toprule
\multicolumn{2}{@{}l}{\textit{Panel A: design contract}} \\
\midrule
Sample timing & Pooled news: \ppvalue{sample-year-start}--\ppvalue{sample-year-end}; complete-case returns: \ppvalue{return-panel-n-observations} common dates, \ppvalue{return-panel-date-start}--\ppvalue{return-panel-date-end}. News and returns are pooled over the same window. \\
Observational units & 100 firms and 4950 unordered dyads $i<j$; one observation per pair and no self-dyads. \\
Metrics & $W_{ij}$ is rooted balanced quadratic Wasserstein distance between row-normalized article clouds under Euclidean chord ground cost; $D_{ij}=\sqrt{2(1-\hat\rho_{ij})}$ from complete-case daily returns. \\
Primary design & Specification 6: $D_{ij}=\mu+a_i+a_j+\beta_{W_2} W_{ij}+u_{ij}$, with symmetric additive firm effects and one firm indicator omitted. \\
Resampling & 1999 multinomial node-bootstrap draws; firms are resampled and dyads receive endpoint-multiplicity weight $w_iw_j$, without artificial self-dyads. \\
Estimand & Same-window contemporaneous association conditional on realized $W$ and $D$ and additive firm margins. \\
\bottomrule
\end{tabularx}
\medskip
\begin{tabularx}{\linewidth}{rrrrrX}
\toprule
\multicolumn{6}{l}{\textit{Panel B: pre-fit conditioning by ladder rung}} \\
\midrule
Spec. & Rank/cols. & Scaled $\kappa$ & VIF$_{W_2}$ & FWL SD ratio & Dropped \\
\midrule
1 & 2/2 & 1.00 & 1.00 & 1.000 & \textemdash \\
2 & 3/3 & 1.81 & 1.39 & 0.847 & \textemdash \\
3 & 4/4 & 1.88 & 1.46 & 0.828 & \textemdash \\
4 & 4/4 & 1.88 & 1.46 & 0.828 & \textemdash \\
5 & 6/6 & 2.31 & 1.84 & 0.738 & \textemdash \\
6 & 101/101 & 13.40 & 2.31 & 0.657 & \textemdash \\
\bottomrule
\end{tabularx}
\caption{Dyadic design and pre-fit audit. Panel A fixes the analysis sample, metrics, specification, resampling unit, and interpretation boundary. Panel B reports conditioning over the complete graph on 100 firms: 4950 dyads, degree 99 for every firm, and 0.040 of unordered dyad pairs share an endpoint. $\kappa$ is computed after centring and unit-norm scaling non-intercept columns. VIF$_{W_2}$ and the FWL residual-SD ratio measure how much W2-distance variation remains after the other regressors and summarize numerical conditioning.}
\label{tab:dyadic-design-diagnostics}
\end{table}


With the point design audited, the remaining inference issue is dependence
among dyads that share a firm.
For bootstrap draw $b$, the $N=\ppvalue{n-firms}$ firms receive counts
\begin{equation}\label{eq:node-bootstrap-counts}
  (w_{1b},\ldots,w_{Nb})
  \sim\operatorname{Multinomial}(N;1/N,\ldots,1/N),
\end{equation}
and dyad $(i,j)$ receives frequency weight $w_{ib}w_{jb}$. The two-sided
sign-tail probability is
\begin{equation}\label{eq:dyadic-bootstrap-tail}
  p_{\mathrm{boot}}
  =2\min\left\{
    \frac{1+\sum_{b=1}^{B}\1\{\widehat\beta_{W_2,b}\leq0\}}{B+1},
    \frac{1+\sum_{b=1}^{B}\1\{\widehat\beta_{W_2,b}\geq0\}}{B+1}
  \right\}\wedge1.
\end{equation}
Resampling firms rather than dyads preserves the shared-endpoint dependence:
when firm $i$ appears $w_{ib}$ times, every incident dyad inherits that
multiplicity through $w_{ib}w_{jb}$.

The add-one correction prevents zero simulated tail probabilities
\citep{phipson2010permutation}. Percentile intervals are first-order summaries
of firm-level sampling uncertainty; they do not incorporate article selection,
encoder estimation, or return measurement as additional generated regressors.

\subsection{Representation and article-sampling sensitivity}
\label{sec:representation-diagnostics}
\label{sec:wasserstein-sensitivity}
\label{sec:xmodel-robustness}
\label{sec:geometry-diagnostics}

Having audited the estimator and inference, we next ask whether the reported
association depends on the particular encoder, width, or distance.
\Cref{tab:w2-representation-sensitivity} holds the estimator, return matrix,
firm sample, and node-bootstrap schedule fixed while changing encoder, width,
or distance.
Rank agreement is Spearman correlation with the canonical Qwen3-Embedding-8B,
128-article W$_2$ matrix; the standardized effect is the change in return
chord distance associated with one dyadic standard deviation of the row's metric.
Writing this effect as $s_r$ for representation $r$, every paired contrast is
$\Delta_{c,r}=s_c-s_r$; a negative value therefore means that the candidate
has the smaller standardized association.
All cell-level association intervals remain positive.
Within these comparisons, W$_1$ nearly preserves the canonical ordering,
while energy produces a materially different ordering but a similar positive
association.
These are sensitivity comparisons, not tests that one representation dominates another.

The lower-cost distance comparators reuse the same deterministic article
samples and Euclidean chord ground metric, so their formulas isolate the
distance functional.
Let $A_{ij}$ denote the cross-sample average chord distance and $V_i,V_j$ the
corresponding within-firm averages.
The reported sample $V$-statistic is
\begin{equation}\label{eq:sample-energy-distance}
  \widehat{\mathcal E}_{ij}
  :=\mathcal E(\widehat Q_i,\widehat Q_j) = \sqrt{
    \frac{2}{N_i N_j} \sum_{n=1}^{N_i} \sum_{m=1}^{N_j}
    d_\Omega(\bm{e}_{i,n}, \bm{e}_{j,m}) - V_i - V_j
  }
\end{equation}
where $d_\Omega(\bm{x}, \bm{y}) = \norm{\bm{x} - \bm{y}}_2$ is the Euclidean
ground metric in $\R^{\ppvalue{embed-dim}}$. Each article embedding is
$L^2$-normalized to unit length, so $d_\Omega$ is the chordal metric on the
unit sphere and $V_i,V_j$ measure within-firm coverage dispersion. The
radicand is nonnegative for the biased $V$-statistic, and the resulting matrix
is symmetric with zero diagonal.

The controlled $W_1$ comparator is the equal-weight balanced assignment
\begin{equation}\label{eq:empirical-wasserstein}
  \widehat W_{1,m}(i,j)
  = \min_{\pi\in S_m}\frac{1}{m}\sum_{a=1}^{m}
  \norm{x_{ia}-x_{j,\pi(a)}}_2.
\end{equation}
It is the minimum average semantic displacement needed to align two firms'
article distributions. Unlike W$_2$, it does not give larger displacements
quadratically greater weight.

The seven-row encoder panel organizes three distinct comparisons.
First, full-width Qwen3-Embedding-8B and 4B compare capacities within one
causal-attention family, while their common 1,024-coordinate outputs isolate
capacity at fixed width.
Second, the 256- and 64-coordinate Qwen3-8B prefixes hold the encoder fixed
and vary its Matryoshka output width.
Third, full-width BGE-large-v1.5 introduces a bidirectional
BERT/RetroMAE-style training lineage from another provider, changing
architecture, pretraining, and provider jointly rather than isolating one
mechanism \citep{xiao_cpack_2023,qwen3_embedding_2025}.
The compact contrast table reports the seven candidate-minus-reference pairs
descriptively with paired intervals and Holm-adjusted $p$-values; none is a
causal or encoder-superiority claim.

% AUTO-GENERATED -- do not edit by hand.
\begin{table}[H]
\centering
\footnotesize
\setlength{\tabcolsep}{3.5pt}
\begin{tabular}{lrrrcr}
\toprule
Encoder & Dim. & Distance & Rank $\rho$ & $\Delta D$/1 SD [95\% CI] & Partial $R^2$ \\
\midrule
\multicolumn{6}{l}{\textit{Panel A: encoder and embedding dimension}} \\
Qwen3-8B & 4096 & $W_2$ & 1.000 & 0.069 [0.056, 0.082] & 0.332 \\
Qwen3-4B & 2560 & $W_2$ & 0.965 & 0.068 [0.055, 0.082] & 0.299 \\
Qwen3-8B & 1024 & $W_2$ & 0.985 & 0.067 [0.055, 0.081] & 0.333 \\
Qwen3-4B & 1024 & $W_2$ & 0.960 & 0.066 [0.054, 0.080] & 0.308 \\
BGE-large-v1.5 & 1024 & $W_2$ & 0.904 & 0.064 [0.050, 0.081] & 0.229 \\
Qwen3-8B & 256 & $W_2$ & 0.925 & 0.062 [0.051, 0.075] & 0.322 \\
Qwen3-8B & 64 & $W_2$ & 0.804 & 0.055 [0.044, 0.066] & 0.279 \\
\midrule
\multicolumn{6}{l}{\textit{Panel B: distance family, canonical encoder and dimension}} \\
Qwen3-8B & 4096 & $W_2$ & 1.000 & 0.069 [0.056, 0.082] & 0.332 \\
Qwen3-8B & 4096 & $W_1$ & 0.999 & 0.069 [0.056, 0.083] & 0.332 \\
Qwen3-8B & 4096 & Energy & 0.289 & 0.085 [0.070, 0.102] & 0.356 \\
\midrule
\multicolumn{6}{l}{\textit{Panel C: matched encoder-vintage negative control}} \\
EttaX V0 & 320 & $W_2$ & 0.747 & 0.053 [0.036, 0.074] & 0.081 \\
EttaX V1 & 320 & $W_2$ & 0.749 & 0.057 [0.041, 0.080] & 0.093 \\
EttaX V3 & 320 & $W_2$ & 0.768 & 0.053 [0.038, 0.073] & 0.093 \\
\bottomrule
\end{tabular}
\medskip
\begin{tabularx}{\linewidth}{Xrrrrrr}
\toprule
Vintage contrast ($c-r$) & Estimate & 95\% CI & Raw $p$ & Holm $p$ & Bound & Equiv. \\
\midrule
EttaX V1 $-$ EttaX V0 & 0.005 & [0.000, 0.010] & 0.056 & 0.168 & 0.012 & Yes \\
EttaX V3 $-$ EttaX V0 & 0.000 & [-0.005, 0.005] & 0.881 & 0.881 & 0.012 & Yes \\
EttaX V3 $-$ EttaX V1 & -0.004 & [-0.010, 0.000] & 0.064 & 0.168 & 0.012 & Yes \\
\bottomrule
\end{tabularx}
\medskip
\begin{tabularx}{\linewidth}{Xrrrrr}
\toprule
Base contrast ($c-r$) & Estimate & 95\% CI & Raw $p$ & Holm $p$ \\
\midrule
Qwen3-4B full $-$ Qwen3-8B full & -0.001 & [-0.003, 0.001] & 0.388 & 1.000 \\
Qwen3-4B@1024 $-$ Qwen3-8B@1024 & -0.001 & [-0.004, 0.002] & 0.437 & 1.000 \\
Qwen3-8B@1024 $-$ Qwen3-8B full & -0.001 & [-0.003, 0.000] & 0.070 & 0.350 \\
Qwen3-8B@256 $-$ Qwen3-8B@1024 & -0.005 & [-0.008, -0.002] & 0.001 & 0.007 \\
Qwen3-8B@64 $-$ Qwen3-8B@256 & -0.008 & [-0.012, -0.004] & 0.001 & 0.007 \\
BGE-large-v1.5 full $-$ Qwen3-8B@1024 & -0.003 & [-0.009, 0.003] & 0.321 & 1.000 \\
BGE-large-v1.5 full $-$ Qwen3-4B@1024 & -0.002 & [-0.008, 0.004] & 0.470 & 1.000 \\
\bottomrule
\end{tabularx}
\caption{Representation sensitivity under the primary estimator. Every row uses symmetric additive firm effects and the same 1,999 node-bootstrap schedule. Rank agreement is Spearman correlation with the canonical Qwen3-Embedding-8B, 128-article rooted-$W_2$ geometry on the common 100-firm, 4,950-dyad panel. The full-width Qwen rows change capacity and native width jointly; the 8B prefixes isolate Matryoshka truncation; BGE-large changes architecture, training pipeline, and provider. The EttaX panel holds architecture, recipe, compute, and sampling fixed while changing only the encoder's Wikipedia snapshot; its paired contrasts reuse identical bootstrap draws and declare equivalence only when the full interval lies inside the configured symmetric bound. Standardized effects use each row's dyadic standard deviation. Raw $p$-values are two-sided add-one sign-tail probabilities; Holm adjustment is applied separately to the seven base and three vintage contrasts. Contrast estimates are candidate-minus-reference differences in $\Delta D$ per one dyadic standard deviation ($c-r$); the vintage equivalence bound is shown in its contrast panel. $W_1$ and energy are lower-cost descriptive comparators; the table does not test representation superiority.}
\label{tab:w2-representation-sensitivity}
\end{table}


Moderate truncation from full width to 1,024 coordinates leaves the 8B
standardized association statistically indistinguishable from the canonical
cell. The subsequent 1,024-to-256 and 256-to-64 compression steps are negative
and remain distinguishable from zero after Holm adjustment. Neither capacity
comparison nor either BGE contrast is distinguishable after the same
base-family adjustment. This pattern locates a compression stress boundary for
the measured association; it does not rank the encoders' general quality.

The base comparisons do not isolate training-data vintage, so Panel C uses a
matched negative control for that narrower dimension.
The three EttaX arms hold architecture, recipe, compute budget, tokenizer,
corpus sampler, and seed fixed while moving the Wikipedia snapshot from
before the sample through a post-sample date.
Their standardized associations remain close to the primary estimate.
Because all cells reuse the same node-bootstrap schedule, the
later-minus-earlier contrasts are paired.
The point contrasts and intervals are reported together with the seven base
contrasts; vintage equivalence retains the symmetric $0.012$ margin
calibrated to the previously computed Qwen3-8B--BGE-large paired contrast.
Within this compact architecture and frozen recipe, all three intervals lie
strictly inside the margin and meet the declared equivalence rule.
The design does not exclude a vintage effect that emerges only at the much
larger capacity of the production encoder.
The positive association is not interpreted as causal, and BGE is not used to
attribute a difference to architecture alone.

Association stability alone does not show how truncation alters the
underlying dyad geometry.
The Matryoshka and deterministic-prefix diagnostics therefore separate
geometry from association.
In \Cref{fig:matryoshka-preservation}, the diagonal is exact distance
preservation and $\rho$ measures dyad-rank agreement with full width.
The 256-coordinate representation is closer to the full matrix than the
64-coordinate representation, although neither prefix is identical.

\begin{figure}[H]
  \centering
  \ppfigure{1}{matryoshka_distance_preservation}
  \caption{Truncated versus full-width Qwen3-Embedding-8B rooted W$_2$ over the
    common firm pairs. Both axes are rooted W$_2$ distance in unit-chord
    embedding space; dashed lines mark equality and $\rho$ is Spearman rank
    agreement with the full-width matrix. The 256-coordinate prefix preserves
    the canonical geometry more closely than the 64-coordinate prefix. These
    projections diagnose truncation, not economic validity. Authors'
  calculations.}
  \label{fig:matryoshka-preservation}
\end{figure}

Compression sensitivity does not address whether the result depends on the
finite article prefix used to construct each firm distribution.
\Cref{tab:w2-sampling-diagnostics} therefore varies the deterministic
per-firm article prefix at fixed encoder and distance.
Standardized effects remain positive and rank agreement increases toward the
128-article benchmark.
The exercise measures sensitivity to the declared sample construction; it is
not a sampling-theory correction for the latent article distribution.

% AUTO-GENERATED -- do not edit by hand.
\begin{table}[H]
\centering
\footnotesize
\setlength{\tabcolsep}{3.5pt}
\begin{tabular}{lrrrcr}
\toprule
Encoder & Articles & Distance & Rank $\rho$ & $\Delta D$/1 SD [95\% CI] & Partial $R^2$ \\
\midrule
Qwen3-8B & 32 & $W_2$ & 0.895 & 0.060 [0.047, 0.073] & 0.260 \\
Qwen3-8B & 64 & $W_2$ & 0.960 & 0.066 [0.054, 0.079] & 0.308 \\
Qwen3-8B & 128 & $W_2$ & 1.000 & 0.069 [0.056, 0.082] & 0.332 \\
\bottomrule
\end{tabular}
\caption{Article-sampling sensitivity for Qwen3-Embedding-8B. Prefixes are deterministic, rank agreement is measured against the 128-article $W_2$ matrix, and all association intervals reuse the common node schedule. The rows diagnose finite article sampling; they do not quantify embedding-estimation uncertainty.}
\label{tab:w2-sampling-diagnostics}
\end{table}


\subsection{Sector structure and projection validity}
\label{sec:projection-validity}
\label{sec:geometric-representation}

The remaining diagnostics separate sector structure in the original distance
geometry from the fidelity of its two-dimensional display in \Cref{fig:sector-mds}.
DISCO (distance components) addresses the first question by decomposing total
energy dispersion in the original distance matrix into within- and
between-sector components \citep{szekely_energy_2013}.
The labels come from Nasdaq summary metadata without a recorded snapshot date.
Sector membership accounts for $R^2_{\mathcal E}=\ppnum[3]{disco-r2}$ of
total dispersion, so sector is relevant but leaves most geometry unexplained.
The permutation statistic asks whether the observed sector grouping is more
distinct than relabelled firms; it does not validate the article proxy or
identify sector as the mechanism behind the W$_2$--return association.

% AUTO-GENERATED by pipeline render-paper1-tables -- do not edit by hand.
\begin{table}[H]
\centering
\small
\textit{Panel A: DISCO decomposition and location}
\par\smallskip
\begin{tabular}{rrrrrr}
\toprule
$S_W$ & $S_B$ & $T$ & $R_E^2$ & $F$ & $p_{\rm perm}$ \\
\midrule
21.000 & 3.783 & 24.783 & 0.153 & 2.049 & 0.0001 \\
\bottomrule
\end{tabular}
\par\medskip
\textit{Panel B: dispersion diagnostic}
\par\smallskip
\begin{tabular}{rr}
\toprule
$F_{\rm disp}$ & $p_{\rm disp}$ \\
\midrule
14.446 & 0.982 \\
\bottomrule
\end{tabular}
\caption{DISCO diagnostics on the canonical energy metric using the cited V-statistic denominators. Panel A decomposes total dispersion and tests sector-location structure. Panel B tests whether sector dispersions differ; it is a diagnostic companion rather than a second location test. Both tests use 9999 sector-label permutations over the fixed firm universe. The principal-coordinate representation has 99 positive axes and minimum eigenvalue \num{9.64e-17}. Sector sizes are Communication Services=11, Consumer Cyclical=18, Consumer Defensive=7, Energy=1, Financial Services=1, Healthcare=14, Industrials=7, Technology=38, Utilities=3.}
\label{tab:disco-decomposition}
\end{table}


The projection audit addresses the second question.
Because the MDS projection (\Cref{fig:sector-mds}) compresses
high-dimensional distances into two coordinates, local-neighbour fidelity is limited.
\Cref{fig:mds-shepard-diagnostic} quantifies this: the projection supports a
global description but not local pair selection.

\begin{figure}[H]
  \centering
  \ppfigure{0.62}{mds_shepard_diagnostic}
  \caption{Fidelity of the two-dimensional metric MDS projection. The
    horizontal axis is exact pairwise W$_2$ and the vertical axis is projected
    Euclidean distance. Stress-1 is \ppnum[3]{mds-stress-1}; Pearson and
    Spearman Shepard correlations are \ppnum[3]{mds-shepard-pearson} and
    \ppnum[3]{mds-shepard-spearman}. Mean top-five-neighbour recall is
    \ppnum[3]{mds-top5-neighbour-recall}, and the exact nearest neighbour is
    preserved for \ppvalue{mds-exact-nearest-preserved} of
    \ppvalue{n-firms} firms. The projection supports a global description but
  not local pair selection. Authors' calculations.}
  \label{fig:mds-shepard-diagnostic}
\end{figure}

% Exploratory Later Outcome Paths moved to main-text §5 opening
% (\Cref{fig:neighbour-outcome-paths}).
\label{sec:exploratory-geometry-outcomes}
% pp-interpretation-hash{paper1-neighbour-outcome-split-v2}{14d0af58abc8f7e7de5a2956c26ef0e2170d3722946055a5ad9b44cae4512926}
